\section{Nilpotent quandle homomorphisms}
\label{section: nilpotent homomorphisms}

In this section, we introduce nilpotent surjective homomorphisms of
quandles and study their structure in terms of coverings.


\subsection{Nilpotent quandles}
\label{subsection: nilpotent quandles}

For a group $G$, the \emph{lower central series} of $G$ is defined inductively by $\grpLCS_0(G)=G$ and $\grpLCS_{i+1}(G)=[G,\grpLCS_i(G)]$. 
The group $G$ is called \emph{nilpotent} if $\grpLCS_c(G)$ is trivial for some $c\geq 0$. 
Remark here that we adopt the degree-zero convention on indexing: the lower central series starts at $0$, while Darn\'e's notation uses indices starting at $1$.

Using the inner automorphism group, Darn\'e introduced nilpotent quandles:

\begin{dfn}[{\cite[Definition~2.1]{darne_2026_nilpotent_quandles}}]
\label[dfn]{definition: nilpotent quandle}
    A quandle $Q$ is called \emph{nilpotent} if $\Inn(Q)$ is a nilpotent group. More precisely, we say that $Q$ is \emph{nilpotent of class at most $c$}, or \emph{$c$-nilpotent}, if $\grpLCS_{c}(\Inn(Q))=1$. The smallest such integer $c$ is called the \emph{nilpotency class} of $Q$ and is denoted by $\ncl(Q)$.
\end{dfn}

\begin{rem}
    As already remarked, our indexing is slightly different from Darn\'e's \cite{darne_2026_nilpotent_quandles}.
    In terms of our notation, his definition of $(c+1)$-nilpotency is $\relLCS_{c}(\Inn(Q))=1$.
    Thus, our numerical class is one less than Darn\'e's quandle-class convention.
\end{rem}

\begin{eg}
    A quandle $Q$ is $0$-nilpotent if and only if $\Inn(Q)$ is a trivial group, if and only if $Q$ is a trivial quandle.
    A quandle $Q$ is $1$-nilpotent if and only if $\Inn(Q)$ is an abelian group.
\end{eg}




\subsection{Nilpotent homomorphisms and their basic properties}
\label{subsection: nilpotent homomorphisms and their basic properties}

In the spirit of relativization, we now introduce a notion of nilpotency for quandle homomorphisms. The nilpotency of the relative inner automorphism group alone, however, is not sufficient; we need to consider the conjugation action of the whole inner automorphism group. In other words, we utilize the notion of relative nilpotency of a normal subgroup.

Recall that for a normal subgroup $N$ of a group $G$, the \emph{$G$-relative lower central series} of $N$ is defined inductively by $\relLCS_0(G,N)=N$ and $\relLCS_{i+1}(G,N)=[G,\relLCS_i(G,N)]$. 
Then, $N$ is called \emph{$G$-nilpotent} if $\relLCS_c(G,N)$ is trivial for some $c\geq 0$. See \cref{subsection in appendix: relative nilpotency for groups} for the details on relative nilpotency of groups.

Applying this to $\relInn(f)\trianglelefteq\Inn(P)$, we introduce the following definition.

\begin{dfn}
\label[dfn]{definition: nilpotent quandle homomorphism}
    A surjective quandle homomorphism $f\colon P\twoheadrightarrow Q$ is called \emph{nilpotent} if $\relInn(f)$ is $\Inn(P)$-nilpotent, or equivalently if the group homomorphism $f_*\colon \Inn(P)\twoheadrightarrow \Inn(Q)$ is nilpotent. More precisely, for $c\geq0$, we say that $f$ is \emph{nilpotent of class at most $c$}, or \emph{$c$-nilpotent}, if $\relLCS_{c}(\Inn(P),\relInn(f))=1$. We write $\ncl(f)\coloneqq \ncl(f_*)=\ncl_{\Inn(P)}(\relInn(f))$ for its nilpotency class.
\end{dfn}

In what follows, we will often use the abbreviation $\relLCS_i(f)\coloneqq \relLCS_i(\Inn(P),\relInn(f))$.

%Note that, in our convention in \cref{appendix: relative nilpotency for groups}, when $f$ is not nilpotent, $\ncl(f)=\infty$.

\begin{eg}
    For a non-empty quandle $P$, consider the unique homomorphism $f\colon P\twoheadrightarrow\{\ast\}$ into the terminal quandle. Then, since $\relInn(f) = \Inn(P)$, $f$ is $c$-nilpotent if and only if $P$ is a $c$-nilpotent quandle.

    On the other hand, we do not consider the empty quandle nilpotent, since the unique homomorphism $\emptyset\to\{\ast\}$ is not surjective. This differs from Darn\'e's convention; see \cref{remark: minor caveat concerning the empty quandle}.
\end{eg}


The following characterizations of the first two nilpotency classes are immediate.

\begin{prop}
\label[prop]{proposition: low nilpotency classes}
    Let $f\colon P\twoheadrightarrow Q$ be a surjective quandle homomorphism.
    \begin{enumerate}
        \item $f$ is $0$-nilpotent if and only if it is rigid.
        \item $f$ is $1$-nilpotent if and only if $f_*\colon \Inn(P)\to \Inn(Q)$ is a central extension.
    \end{enumerate}
\end{prop}

In particular, we have the implications: $0$-nilpotent $\Leftrightarrow$ rigid $\Rightarrow$ covering $\Rightarrow$ $1$-nilpotent.


\begin{prop}
\label[prop]{proposition: nilpotency under composition}
Let $f\colon P\twoheadrightarrow Q$ and $g\colon Q\twoheadrightarrow R$ be surjective quandle homomorphisms. 
\begin{enumerate}
    \item\label{item:composition_of_nilpotent_quandle_hom} If $f$ is $c$-nilpotent and $g$ is $d$-nilpotent, then $g \circ f$ is $(c+d)$-nilpotent.
    \item\label{item:cancelation_of_nilpotent_quandle_hom} If $g\circ f$ is $c$-nilpotent, then both $f$ and $g$ are so.
\end{enumerate}
Consequently, whenever all three homomorphisms are nilpotent,
\[
\max\{\ncl(f),\ncl(g)\}\leq\ncl(g\circ f) \leq\ncl(f)+\ncl(g).
\]
\end{prop}

\begin{proof}
By the functoriality of $\Inn(-)$ for surjections (\cref{proposition: functoriality of Inn for surjections}), we have surjective group homomorphisms $f_*\colon \Inn(P)\to \Inn(Q)$ and $g_*\colon \Inn(Q)\to\Inn(R)$.
Applying \cref{proposition: composition of nilpotent group homomorphisms} to these homomorphisms yields the assertions.
\end{proof}

\begin{prop}[cf. {\cite[Proposition 2.9]{ishihara_tamaru_2016_flat_connected_finite_quandles}}]
\label[prop]{proposition: inner automorphism groups of products}
Let $P$ and $P'$ be quandles. Then, there is a canonical injective group homomorphism
\[ \iota_{P,P'}\colon \Inn(P\times P') \hookrightarrow \Inn(P)\times\Inn(P') \]
satisfying $\iota_{P,P'}(s_{(x,x')})=(s_x,s_{x'})$. Moreover, these homomorphisms are natural with respect to surjective quandle homomorphisms.
\end{prop}

\begin{proof}
If either factor is empty, then $\Inn(P\times P')=1$, and the assertions are immediate.
Suppose that $P$ and $P'$ are non-empty, and let $\pi_1\colon P\times P'\twoheadrightarrow P$ and $\pi_2\colon P\times P'\twoheadrightarrow P'$ be the projections. Then, we obtain a natural group homomorphism $\iota_{P,P'}\coloneqq((\pi_1)_*,(\pi_2)_*)\colon \Inn(P\times P')\to \Inn(P)\times\Inn(P')$ from the universal property of the product of groups.

Let $\varphi\in\Ker(\iota_{P,P'})$. Then, $(\pi_1)_*(\varphi)=\id=(\pi_2)_*(\varphi)$, so $\pi_1\circ\varphi=\pi_1$ and $\pi_2\circ\varphi=\pi_2$. Thus, the universal property of $P\times P'$ gives $\varphi=\id_{P\times P'}$. This shows $\iota_{P,P'}$ is injective.

For surjective homomorphisms $f\colon P\twoheadrightarrow Q$ and $f'\colon P'\twoheadrightarrow Q'$, both $(f_*\times f'_*)\circ\iota_{P,P'}$ and $\iota_{Q,Q'}\circ(f\times f')_*$ send $s_{(x,x')}$ to $(s_{f(x)},s_{f'(x')})$. Since the symmetry maps generate $\Inn(P\times P')$, these homomorphisms coincide, proving naturality.
\end{proof}

\begin{prop}
\label[prop]{proposition: nilpotency under finite products}
    Let $f\colon P\twoheadrightarrow Q$ and $f'\colon P'\twoheadrightarrow Q'$ be $c$- and $c'$-nilpotent, respectively. Then $f\times f'\colon P\times P'\twoheadrightarrow Q\times Q'$ is nilpotent of class at most $\max\{c,c'\}$.
\end{prop}

\begin{proof}
By \cref{proposition: inner automorphism groups of products}, we have an inclusion between surjections of groups:
\[\begin{tikzcd}
    \Inn(P)\times \Inn(P') \arrow[two heads]{r}{f_*\times f'_*} & \Inn(Q) \times \Inn(Q') \\
    \Inn(P\times P') \arrow[two heads]{r}{(f\times f')_*} \arrow[hook]{u} & \Inn(Q\times Q')\rlap{.} \arrow[hook]{u}
\end{tikzcd}\]
Since $f_*$ and $f'_*$ are nilpotent of class at most $c$ and $c'$ by the assumption, respectively, the product $f_*\times f'_*$ is nilpotent of class at most $d\coloneqq\max\{c,c'\}$ by \cref{proposition: relative nilpotency under finite products}. Therefore, by \cref{lemma: relative lower central series under subgroups}, so is $(f\times f')_*$, and thus the product $f\times f'$ is $d$-nilpotent.
\end{proof}


The relative-inner filtration records the number of covering steps before the final rigid factor.
The characterization below is a direct relativization of \cite[Theorem 2.13]{darne_2026_nilpotent_quandles}.

\begin{lem}
\label[lem]{lemma: quotient covering via commutator}
    Let $P$ be a quandle, and $N\subseteq M\subseteq \Inn(P)$ be (normal) subgroups of $\Inn(P)$. If $[\Inn(P),M] \subseteq N$, then the induced quandle homomorphism $h \colon P/N \to P/M$ is a covering homomorphism.
\end{lem}
\begin{proof}
%Since $[\Inn(P),N]\subseteq[\Inn(P),M]\subseteq N$ and $[\Inn(P),M]\subseteq N\subseteq M$, both $N$ and $M$ are normal in $\Inn(P)$. Thus, the quotients $P/N$ and $P/M$ are defined.
Let $\pi_N$ and $\pi_M$ be the quotient map by $N$ and $M$, respectively. Take $\pi_N(x),\pi_N(y)\in P/N$, where $x,y \in P$, such that they have the same image in $P/M$. Then, $\pi_M(x)=h(\pi_N(x))=h(\pi_N(y))=\pi_M(y)$, so we have $y=\varphi(x)$ for some $\varphi\in M$. The assumption implies that
\[ s_xs_y^{-1} =s_xs_{\varphi(x)}^{-1} = s_x\varphi s_x^{-1} \varphi^{-1} =[s_x,\varphi] \in [\Inn(P),M] \subseteq N. \]
Thus, it follows that $(\pi_N)_*(s_xs_y^{-1})=\id$, that is, $s_{\pi_N(x)}=s_{\pi_N(y)}$ on $P/N$. Therefore, $P/N\twoheadrightarrow P/M$ is a covering homomorphism.
\end{proof}

\begin{thm}
\label[thm]{theorem: nilpotent homomorphism and covering rigid tower}
Let $f\colon P\twoheadrightarrow Q$ be a surjective quandle homomorphism and let $c\geq0$. Then the following conditions are equivalent.
\begin{enumerate}
    \item\label{item:c-nilpotent} $f$ is $c$-nilpotent.
    \item\label{item:factorization_into_c_covering_and_1_rigid} $f$ admits a factorization
    \[
    \begin{tikzcd}
        P=P_{c}
        \arrow[two heads]{r}{p_c}
        &
        P_{c-1}
        \arrow[two heads]{r}{p_{c-1}}
        &
        \cdots
        \arrow[two heads]{r}
        &
        P_{1}
        \arrow[two heads]{r}{p_1}
        &
        P_{0}
        \arrow[two heads]{r}{h}
        &
        Q
    \end{tikzcd}
    \]
    in which the first $c$ homomorphisms $p_c,\dots,p_1$ are coverings and the last homomorphism $h$ is rigid.
\end{enumerate}
\end{thm}

\begin{proof}
\eqref{item:c-nilpotent} $\Rightarrow$ \eqref{item:factorization_into_c_covering_and_1_rigid}:
We abbreviate $\relLCS_i(\Inn(P),\relInn(f))$ to $\relLCS_i(f)$. By \cref{theorem: connected-rigid factorization}, $f$ factors as a connected homomorphism $p\colon P\to P/\relInn(f)$ followed by a rigid homomorphism $h\colon P/\relInn(f) \to Q$. If $f$ is $c$-nilpotent, then $\relLCS_c(f)=1$, and hence the $\Inn(P)$-relative lower central series
\[ 1= \relLCS_c(f) \subseteq \relLCS_{c-1}(f) \subseteq \cdots \subseteq \relLCS_0(f) = \relInn(f) \]
stops. Taking successive quotients by these normal subgroups, we moreover obtain the following factorization of $p$:
\[
\begin{tikzcd}
    P=P/\relLCS_c(f)
    \arrow[two heads]{r}{p_c}
    &
    P/\relLCS_{c-1}(f)
    \arrow[two heads]{r}{p_{c-1}}
    &
    \cdots
    \arrow[two heads]{r}{p_1}
    &
    P/\relLCS_0(f) = P/\relInn(f).
\end{tikzcd}
\]
By \cref{lemma: quotient covering via commutator}, all $p_i$ are in fact a covering homomorphism. Thus $f$ has the required factorization.

\eqref{item:factorization_into_c_covering_and_1_rigid} $\Rightarrow$ \eqref{item:c-nilpotent}: By \cref{proposition: covering hom is central extension}, applying $\Inn(-)$ to the factorization of $f$ induces a factorization of $f_*$ into $c$ central extensions. Therefore, by \cref{theorem: characterization of relative nilpotency}, the group homomorphism $f_*$ is nilpotent of class at most $c$, and hence $f$ is $c$-nilpotent.
\end{proof}

By \cref{theorem: connected-rigid factorization}, the rigid part of the connected--rigid factorization of a connected homomorphism is an isomorphism. Thus, we have the following corollary.

\begin{cor}
\label[cor]{corollary: connected nilpotent homomorphism}
If $f\colon P\twoheadrightarrow Q$ is connected, then $f$ is $c$-nilpotent if and only if it is a composite of $c$ coverings.
\end{cor}

\begin{prop}
    Let $f\colon P\to Q$ be a surjective quandle homomorphism and $k\colon Q'\to Q$ be an arbitrary homomorphism. Consider the pullback $f'\colon P'\coloneqq P\times_Q Q'\to Q'$ of $f$ along $k$. If $f$ is $c$-nilpotent, then the pullback $f'$ is $(c+1)$-nilpotent.
\end{prop}

\begin{proof}
    If $f$ is $c$-nilpotent, then $f$ is a composition of $c+1$ covering homomorphisms by \cref{theorem: nilpotent homomorphism and covering rigid tower}, since rigid homomorphisms are covering. Pulling back this chain yields a factorization of $f'$ into $c+1$ covering homomorphisms, by the stability of covering homomorphisms under pullback (\cref{proposition: pull-back of covering is again covering}). Applying $\Inn(-)$ to this factorization, we obtain a decomposition of $f'_*$ into $c+1$ central extensions. Hence, $f'$ is $(c+1)$-nilpotent by \cref{theorem: characterization of relative nilpotency}.
\end{proof}

\begin{cor}
\label[cor]{corollary: fibers of nilpotent homomorphisms}
If $f\colon P\twoheadrightarrow Q$ is $c$-nilpotent, then every fiber
$f^{-1}(q)$, regarded as a subquandle of $P$, is $c$-nilpotent.
\end{cor}

\begin{proof}
   Consider the pullback of the factorization in \cref{theorem: nilpotent homomorphism and covering rigid tower} along $\{q\}\hookrightarrow Q$. The first $c$ homomorphisms remain coverings, and the last one is rigid since its domain is a fiber of a covering homomorphism and hence a trivial quandle. Thus the assertion follows from \cref{theorem: nilpotent homomorphism and covering rigid tower} again.
\end{proof}



Darn\'e's universal nilpotent quotient construction in \cite[Proposition~2.15]{darne_2026_nilpotent_quandles} admits the following relative version.

\begin{prop}
\label[prop]{proposition: universal property of nilpotent truncation}
    Let $f\colon P\to Q$ be a surjective quandle homomorphism and let $c\geq0$. Then, the induced homomorphism $\nilpTrunc{c}{f}\colon P/\relInnLCS{c}{f} \twoheadrightarrow Q$ is $c$-nilpotent. 
    
    Moreover, the factorization $f= \nilpTrunc{c}{f}\circ \pi$, where $\pi\colon P\to P/\relInnLCS{c}{f}$ is the quotient map, is universal among factorizations of $f$ whose second homomorphism is $c$-nilpotent.
    That is, if $f=h\circ u$ with $u\colon P\twoheadrightarrow R$ surjective and $h\colon R\twoheadrightarrow Q$ $c$-nilpotent, then $u$ factors uniquely through $P/\relInnLCS{c}{f}$.
\end{prop}

\begin{proof}
    Note that since $\relInnLCS{c}{f}\subseteq \relInn(f)$, $f$ factors through the quotient map $\pi\colon P\to P/\relInnLCS{c}{f}$ as $f=\nilpTrunc{c}{f}\circ \pi$ (\cref{prop:basic_property_of_quotient_by_normal_subgroup}). Applying $\Inn(-)$, we have $f_*= (\nilpTrunc{c}{f})_*\circ \pi_*$. Since we see from \cref{lemma: relative symmetry groups under a surjective factorization}
    \[\pi_*(\relInn(f))=\pi_*(\Ker(f_*))= \Ker((\nilpTrunc{c}{f})_*)= \relInn(\nilpTrunc{c}{f}), \]
    \cref{proposition: functoriality of relative lower central series} implies that
    \[ \relInnLCS{c}{\nilpTrunc{c}{f}} = \relLCS_c(\Inn(P/\relInnLCS{c}{f}),\relInn(\nilpTrunc{c}{f})) = \pi_*( \relLCS_c(\Inn(P),\relInn(f)))=\pi_*(\relInnLCS{c}{f}) =1, \]
    which shows that $\nilpTrunc{c}{f}$ is $c$-nilpotent.
    
    For the universal property, if $f=h \circ u$ and $h$ is $c$-nilpotent, then $u_*(\relInnLCS{c}{f})=\relInnLCS{c}{h}=1$, so $u$ factors uniquely through $P/\relInnLCS{c}{f}$ by \cref{prop:basic_property_of_quotient_by_normal_subgroup}.
\end{proof}

\begin{dfn}
\label[dfn]{definition: nilpotent truncations}
    Let $c\geq0$. For a surjective quandle homomorphism $f\colon P \twoheadrightarrow Q$, the induced $c$-nilpotent homomorphism $\nilpTrunc{c}{f}\colon P/\relInnLCS{c}{f}\twoheadrightarrow Q$ is called the \emph{$c$-nilpotent truncation} of $f$.
\end{dfn}



\subsection{Nilpotency via relative transvection groups}
\label{subsection: Nilpotency via relative transvections and adjoint groups}


We have defined the nilpotency of $f$ in terms of the relative inner group $\relInn(f)$. We next show that the same nilpotency condition can also be characterized in terms of the relative transvection group $\relTrans(f)$. The following inclusion provides the key link between these two groups.

\begin{lem}[{\cite[Lemma 3.11]{preprint_yuki_tomoki_2026_relativization_of_symmetries_on_quandles}}]
\label{lem:commutator_inn_rel_trans_rel}
    For a surjective quandle homomorphism $f \colon P \twoheadrightarrow Q$, we have $[\Inn(P), \relInn(f)] \subseteq \relTrans(f)$.
    Moreover, if $f$ is connected, the equality $[\Inn(P), \relInn(f)] = \relTrans(f)$ holds.
\end{lem}


We have defined $\relInnLCS{i}{f}\coloneqq \relLCS_i(\Inn(P),\relInn(f))$ for a surjective homomorphism $f\colon P \twoheadrightarrow Q$. We also define $\relTransLCS{i}{f}\coloneqq\relLCS_{i}(\Inn(P),\relTrans(f))$ for the $i$-th term of the relative lower central series of $\relTrans(f)$.



\begin{prop}
\label[prop]{proposition: nilpotency detected by relative transvection group}
Let $f\colon P\twoheadrightarrow Q$ be a surjective homomorphism, and let $c\geq 0$.
\begin{enumerate}
    \item\label{item:nilpotent_of_relInn_implies_that_of_relTrans} If $\relInn(f)$ is $\Inn(P)$-nilpotent of class at most $c$, then $\relTrans(f)$ is $\Inn(P)$-nilpotent of class at most $c$.
    \item\label{item:nilpotent_of_relTrans_implies_that_of_relInn} If $\relTrans(f)$ is $\Inn(P)$-nilpotent of class at most $c$, then $\relInn(f)$ is $\Inn(P)$-nilpotent of class at most $c+1$.
\end{enumerate}
\end{prop}

\begin{proof}
\eqref{item:nilpotent_of_relInn_implies_that_of_relTrans}: For all $i\geq 0$, the inclusions $\relTransLCS{i}{f}\subseteq \relInnLCS{i}{f}$ follow inductively from $\relTrans(f)\subseteq \relInn(f)$. Therefore, if $\relInnLCS{c}{f}=1$, then we have $\relTransLCS{c}{f}=1$.

\eqref{item:nilpotent_of_relTrans_implies_that_of_relInn}: \cref{lem:commutator_inn_rel_trans_rel} says $\relInnLCS{1}{f}=[\Inn(P),\relInn(f)]\subseteq \relTrans(f)=\relTransLCS{0}{f}$. Thus, by induction again, we see $\relInnLCS{i+1}{f}=\relLCS_i(\Inn(P),\relInnLCS{1}{f}) \subseteq \relLCS_i(\Inn(P),\relTransLCS{0}{f})=\relTransLCS{i}{f}$.
From this, if $\relTransLCS{c}{f}=1$, then we have $\relInnLCS{c+1}{f}=1$.
\end{proof}

From \cite[Corollary 3.7]{darne_2026_nilpotent_quandles}, any connected and nilpotent quandle is isomorphic to the terminal quandle. This can be relativized as follows.

\begin{prop}
\label[prop]{proposition: strongly connected and nilpotent is iso}
    Any strongly connected and nilpotent homomorphism is an isomorphism.
\end{prop}

\begin{proof}
    If $f\colon P \twoheadrightarrow Q$ is strongly connected, then $\relTrans(f)$ is $\Inn(P)$-perfect by \cref{proposition:relTrans_of_strongly_conncted_is_Inn-perfect}, and $\relTransLCS{i}{f}=\relTrans(f)$ for all $i\geq 0$. Thus, if $f$ is moreover $c$-nilpotent for some $c\geq 0$, then $\relInnLCS{c}{f}$, hence $\relTransLCS{c}{f}=\relTrans(f)$, becomes trivial.
    Since $f$ is isomorphic to $\pi_{\relTrans(f)}$ as quotients of $P$, $f$ is an isomorphism.
\end{proof}


The nilpotency class of $\relTrans(f)$, plus one, turns out to record the length of a decomposition into covering homomorphisms.


\begin{thm}
\label[thm]{theorem: transvection class and iterated coverings}
Let $f\colon P\twoheadrightarrow Q$ be a surjective quandle homomorphism and let $c\geq0$. Then, $\relTransLCS{c}{f}=1$ if and only if $f$ is a composition of $c+1$ covering homomorphisms.
\end{thm}


\begin{proof}
($\Rightarrow$): By \cref{theorem: maximal covering factorization}, $f$ factors as a homomorphism $p\colon P\to P/\relTrans(f)$ followed by a covering homomorphism $h\colon P/\relTrans(f) \to Q$.
If $\relTransLCS{c}{f}=1$, then the $\Inn(P)$-relative lower central series
\[ 1= \relTransLCS{c}{f} \subseteq \relTransLCS{c-1}{f} \subseteq \cdots \subseteq \relTransLCS{0}{f} = \relTrans(f) \]
stops. Taking successive quotients by these normal subgroups, we moreover obtain the following factorization of $p$:
\[
\begin{tikzcd}
    P=P/\relTransLCS{c}{f}
    \arrow[two heads]{r}{p_c}
    &
    P/\relTransLCS{c-1}{f}
    \arrow[two heads]{r}{p_{c-1}}
    &
    \cdots
    \arrow[two heads]{r}{p_1}
    &
    P/\relTransLCS{0}{f}= P/\relTrans(f).
\end{tikzcd}
\]
By \cref{lemma: quotient covering via commutator}, all $p_i$ are in fact covering homomorphisms. Thus $f$ has the desired factorization of $(c+1)$ covering homomorphisms.

($\Leftarrow$): We show by induction on $c$. The case $c=0$ is the characterization of covering homomorphisms (see \cref{lemma: elementary properties of relative symmetry groups}). Next suppose that $f=h\circ p$, where $p\colon P\twoheadrightarrow R$ is a covering and $h\colon R\twoheadrightarrow Q$ is a composition of $c$ coverings. By \cref{lemma: relative symmetry groups under a surjective factorization}, the map $p_*$ sends $\relTrans(f)$ onto $\relTrans(h)$. By the induction hypothesis, $\relTransLCS{c-1}{h}=\relLCS_{c-1}(\Inn(R),\relTrans(h))=1$, so \cref{proposition: functoriality of relative lower central series} gives $p_*(\relTransLCS{c-1}{f})=\relTransLCS{c-1}{h}=1$, and hence $\relTransLCS{c-1}{f}\subseteq\Ker(p_*)$. Since $\Ker(p_*)\subseteq \Inn(P)$ is central by \cref{proposition: covering hom is central extension}, we thus obtain $\relTransLCS{c}{f}=[\Inn(P),\relTransLCS{c-1}{f}]=1$.
\end{proof}


\begin{dfn}
\label[dfn]{definition: covering length and transvection nilpotency class}
    Let $f\colon P\to Q$ be a surjective quandle homomorphism.
    We define the \emph{covering length} $\covlen(f)$ as the smallest integer $n\geq 1$ for which $f$ can be expressed as a composition of $n$ non-invertible covering homomorphisms.
    By convention, isomorphisms are understood to have covering length $0$.
    If $f$ does not admit a covering decomposition, set $\covlen(f)=\infty$.
\end{dfn}

In addition, we put $\tncl(f)\coloneqq\ncl_{\Inn(P)}(\relTrans(f))$.
The above \cref{theorem: transvection class and iterated coverings} says $\tncl(f)+1 = \covlen(f)$ for a non-isomorphism $f$.


\begin{prop}
\label[prop]{theorem: universal property of covering length truncation}
    Let $f\colon P\to Q$ be a surjective quandle homomorphism and let $c\geq 0$. Then, the induced homomorphism $\covTrunc{c}{f}\colon P/\relTransLCS{c}{f} \twoheadrightarrow Q$ admits a decomposition into covering homomorphisms of length at most $c+1$. 
    
    Moreover, the factorization $f= \covTrunc{c}{f}\circ \pi$, where $\pi\colon P\to P/\relTransLCS{c}{f}$ is the quotient map, is universal among factorizations of $f$ whose second homomorphism has covering length at most $c+1$.
    That is, if $f=h\circ u$ with $u\colon P\twoheadrightarrow R$ surjective and $h\colon R\twoheadrightarrow Q$ of covering length at most $c+1$, then $u$ factors uniquely through $P/\relTransLCS{c}{f}$.
\end{prop}

\begin{proof}
    Note that since $\relTransLCS{c}{f}\subseteq \relInn(f)$, $f$ factors through the quotient map $\pi\colon P\to P/\relTransLCS{c}{f}$ as $f=\covTrunc{c}{f}\circ \pi$ (\cref{prop:basic_property_of_quotient_by_normal_subgroup}). 
    From \cref{lemma: relative symmetry groups under a surjective factorization} and \cref{proposition: functoriality of relative lower central series}, it follows that $\pi_*(\relTransLCS{c}{f})=\relTransLCS{c}{\covTrunc{c}{f}}$. Since $\pi_*(\relTransLCS{c}{f})=1$, we have $\relTransLCS{c}{\covTrunc{c}{f}}=1$, and hence $\covTrunc{c}{f}$ is a composition of $c+1$ covering homomorphisms by \cref{theorem: transvection class and iterated coverings}.
    
    For the universal property, if $f=h \circ u$ and $h$ is a composition of at most $c+1$ coverings, then $\relTransLCS{c}{h}=1$ by \cref{theorem: transvection class and iterated coverings}. Hence, $u_*(\relTransLCS{c}{f})=\relTransLCS{c}{h}=1$, and $u$ factors uniquely through $P/\relTransLCS{c}{f}$ by \cref{prop:basic_property_of_quotient_by_normal_subgroup}.
\end{proof}



\begin{dfn}
\label[dfn]{definition: covering length truncations}
    Let $c\geq 1$. For a surjective quandle homomorphism $f\colon P \twoheadrightarrow Q$, the induced homomorphism $\covTrunc{c-1}{f}\colon P/\relTransLCS{c-1}{f}\twoheadrightarrow Q$ of covering length at most $c$ is called the \emph{$c$-th covering-length truncation} of $f$.
\end{dfn}




\subsection{Nilpotency via relative adjoint groups}
\label{subsection: Nilpotency via relative adjoint groups}


The relative adjoint group, defined below, gives a third group-theoretic realization of the nilpotency condition. While the two groups $\relInn(f)$ and $\relTrans(f)$ are subgroups of $\Inn(P)$, the relative adjoint group lives in the central extension $\As(P)\twoheadrightarrow\Inn(P)$ and therefore changes the class by at most one. This is the relative counterpart of \cite[Corollary~2.5]{darne_2026_nilpotent_quandles}.

\begin{dfn}
\label[dfn]{definition: relative adjoint group}
For a surjective quandle homomorphism $f\colon P\twoheadrightarrow Q$, we define
\[
\relAs(f)\coloneqq\Ker(\As(f)\colon \As(P)\twoheadrightarrow\As(Q))
\]
and call it the \emph{relative adjoint group} of $f$.
\end{dfn}

\begin{lem}
\label[lem]{lemma: generators of relative adjoint group}
For a surjective homomorphism $f\colon P\twoheadrightarrow Q$, we have
\[
    \relAs(f) = \left\langle g_{x}g_{y}^{-1} \mid x,y\in P \text{ such that } f(x)=f(y) \right\rangle.
\]
\end{lem}

\begin{proof}
Let $E$ denote the subgroup on the right. Since $\As(f)(g_{x}g_{y}^{-1})=1$ whenever $f(x)=f(y)$, we have $E\subseteq\relAs(f)$.

First note that $E$ is normal in $\As(P)$. Indeed, $g_{z}(g_{x}g_{y}^{-1})g_{z}^{-1}=g_{s_z(x)}g_{s_z(y)}^{-1}$, and $f(s_z(x))=f(s_z(y))$ whenever $f(x)=f(y)$. Since $s_z$ is bijective, conjugation by $g_{z}$ permutes the generators of $E$.

Let $\pi\colon\As(P)\twoheadrightarrow\As(P)/E$ be the quotient map. For $q\in Q$, define $\overline{g}_{q}$ to be the common value $\pi(g_{x})$ for $x\in f^{-1}(q)$. If $x,x'\in f^{-1}(q)$, then $g_xg_{x'}^{-1}\in E$, and hence $\pi(g_x)=\pi(g_{x'})$. Thus, $\overline{g}_q $ is well-defined. If $f(x)=q$ and $f(y)=r$, then $\overline{g}_{q}\overline{g}_{r} =\pi(g_{x}g_{y}) =\pi(g_{s_x(y)}g_{x}) =\overline{g}_{s_q(r)}\overline{g}_{q}$.
Hence the assignment $g_{q}\mapsto\overline{g}_{q}$ induces a homomorphism $\theta\colon\As(Q)\to\As(P)/E$. For every $x\in P$, $(\theta\circ\As(f))(g_{x})=\pi(g_{x})$, so $\theta\circ\As(f)=\pi$. If $a\in\relAs(f)$, then $\pi(a)=\theta(\As(f)(a))=1$, whence $a\in E$. Thus $\relAs(f)\subseteq E$.
\end{proof}

\begin{prop}
\label[prop]{proposition: relative adjoint group central extension}
    The restriction of $\rho_{P}\colon\As(P)\twoheadrightarrow\Inn(P)$ induces a central extension $\rho_{f}\colon\relAs(f)\twoheadrightarrow\relTrans(f)$.
\end{prop}

\begin{proof}
We see from \cref{lemma: generators of relative adjoint group} that $\rho_{P}(\relAs(f))=\relTrans(f)$, so $\rho_{P}$ restricts to a surjection $\rho_{f}\colon \relAs(f)\twoheadrightarrow \relTrans(f)$. Its kernel $\Ker(\rho_f)= \Ker(\rho_{P}) \cap \relAs(f)$ is contained in $Z(\As(P))\cap \relAs(f) \subseteq Z(\relAs(f))$, since $\Ker(\rho_P) \subseteq Z(\Adj(P))$ by \cref{proposition: elementary properties of adjoint groups}.
\end{proof}

\begin{prop}
\label[prop]{proposition: nilpotency detected by relative adjoint group}
Let $f\colon P\twoheadrightarrow Q$ be a surjective homomorphism, and let $c\geq 0$.
\begin{enumerate}
    \item\label{item:nilpotent_of_relAs_implies_that_of_relTrans} If $\relAs(f)$ is $\As(P)$-nilpotent of class at most $c$, then $\relTrans(f)$ is $\Inn(P)$-nilpotent of class at most $c$.
    \item\label{item:nilpotent_of_relTrans_implies_that_of_relAs} If $\relTrans(f)$ is $\Inn(P)$-nilpotent of class at most $c$, then $\relAs(f)$ is $\As(P)$-nilpotent of class at most $c+1$.
\end{enumerate}
\end{prop}

\begin{proof}
    By \cref{lemma: generators of relative adjoint group}, the surjection $\rho_{P}$ maps $\relAs(f)$ onto $\relTrans(f)$. Hence it follows from \cref{proposition: functoriality of relative lower central series} that $\rho_P(\relLCS_i(\As(P),\relAs(f))) = \relLCS_i(\Inn(P),\relTrans(f))=\relTransLCS{i}{f}$ for every $i\geq 0$.

    \eqref{item:nilpotent_of_relAs_implies_that_of_relTrans}: Thus, if $\relLCS_c(\As(P),\relAs(f))$ is trivial, then so is $\relTransLCS{c}{f}$.

    \eqref{item:nilpotent_of_relTrans_implies_that_of_relAs}: If $\relTransLCS{c}{f}=1$, then $\relLCS_{c}(\As(P), \relAs(f))\subseteq\Ker(\rho_{P})$. The latter subgroup is central in $\As(P)$, so we have $\relLCS_{c+1}(\As(P), \relAs(f))=1$.
\end{proof}



Summarizing the above results, we arrive at the following theorem.

\begin{thm}
\label[thm]{theorem: main structure theorem for nilpotent homomorphisms}
Let $f\colon P\twoheadrightarrow Q$ be a surjective quandle
homomorphism. Then the following conditions are equivalent.
\begin{enumerate}[label={(\roman*)}]
    \item\label{item: f is nilp} $f$ is nilpotent, that is, $\relInn(f)$ is $\Inn(P)$-nilpotent.
    \item\label{item: reltrans is nilp} $\relTrans(f)$ is $\Inn(P)$-nilpotent.
    \item\label{item: relas is nilp} $\relAs(f)$ is $\As(P)$-nilpotent.
    \item\label{item: composition of coverings} $f$ is a finite composition of covering homomorphisms.
\end{enumerate}
\end{thm}

\begin{proof}
The equivalence of \ref{item: f is nilp} and \ref{item: reltrans is nilp} follows from \cref{proposition: nilpotency detected by relative transvection group},
the equivalence of \ref{item: reltrans is nilp} and \ref{item: relas is nilp} follows from \cref{proposition: nilpotency detected by relative adjoint group}, and
the equivalence of \ref{item: reltrans is nilp} and \ref{item: composition of coverings} follows from \cref{theorem: transvection class and iterated coverings}.
\end{proof}



\begin{cor}
\label{corollary: characterization_of_nilpitent_quandles}
    For a non-empty quandle $P$, the following are equivalent.
    \begin{enumerate}[label={(\roman*)}]
    \item\label{item: P is nilp} $P$ is nilpotent, that is, $\Inn(P)$ is nilpotent.
    \item\label{item: Trans is nilp} $\Trans(P)$ is $\Inn(P)$-nilpotent.
    \item\label{item: Adj_zero is nilp} $\As_0(P)=\relAs(P\to \{\ast\})$ is $\As(P)$-nilpotent.
    \item\label{item: Adj is nilp} $\As(P)$ is nilpotent.
    \item\label{item: constructed by iterated coverings} The map $P\to \{\ast\}$ is a finite composition of covering homomorphisms.
\end{enumerate}
\end{cor}

\begin{proof}
\ref{item: Adj_zero is nilp} $\Leftrightarrow$ \ref{item: Adj is nilp}: 
Since $\As(\{*\})\cong\ZZ$ and $\As(P)\to\As(\{*\})$ is the degree homomorphism, we have $\relAs(P\to\{*\})=\As_0(P)$ and $\As(P)/\As_0(P)\cong\ZZ$. 
Thus, \ref{item: Adj_zero is nilp} and \ref{item: Adj is nilp} are equivalent by \cref{lemma: lower central series with cyclic quotient}.

The other equivalences follow from \cref{theorem: main structure theorem for nilpotent homomorphisms}.
\end{proof}


\begin{rem}
    The equivalences among \ref{item: P is nilp}, \ref{item: Adj is nilp}, and \ref{item: constructed by iterated coverings} in \cref{corollary: characterization_of_nilpitent_quandles} were previously established in \cite[Corollary~2.5 and Theorem~2.13]{darne_2026_nilpotent_quandles}, while the equivalences with \ref{item: Trans is nilp} and \ref{item: Adj_zero is nilp} are new.
\end{rem}

\begin{rem}
\label{remark: minor caveat concerning the empty quandle}
    There is one minor point to note concerning the empty quandle. In Darn\'{e}'s sense, a quandle $Q$ is nilpotent if its inner automorphism group $\Inn(Q)$ is nilpotent. Thus, the empty quandle is nilpotent in this sense, since $\Inn(\emptyset)=\{\id\}$. On the other hand, the unique homomorphism $\emptyset\to\{\ast\}$ is not surjective and hence is not a finite composition of covering homomorphisms. Thus, the two notions of nilpotency differ only in the case of the empty quandle.
\end{rem}

\begin{rem}
    Nilpotency for quandles is also studied by Bonatto and Stanovsk\'{y}~\cite{bonatto_stanovsky_2021_commutator_theory_for_racks_and_quandles} from the universal-algebraic viewpoint. According to \cite[Theorem 1.2]{bonatto_stanovsky_2021_commutator_theory_for_racks_and_quandles}, the nilpotency of a quandle $P$ in their sense is equivalent to that of $\Trans(P)$. Hence, our notion of nilpotency is stronger than theirs, as noted by Darn\'{e} in the introduction of \cite{darne_2026_nilpotent_quandles}.
\end{rem}


 
\subsection{Comparison of nilpotency classes and covering length}
\label{subsection: Comparison of nilpotency classes and covering length}

We summarize and compare the three numerical invariants: the $\Inn(P)$-nilpotency class $\ncl(f)$ of $\relInn(f)$, the $\Inn(P)$-nilpotency class $\tncl(f)$ of $\relTrans(f)$, and the covering length $\covlen(f)$.

\begin{thm}
\label{theorem: comparison of ncl tncl covlen}
Let $f\colon P\twoheadrightarrow Q$ be a surjective quandle homomorphism. Then,
\begin{align*}
    \tncl(f)&\leq \ncl(f)\leq \tncl(f)+1,\text{ and} \\
    \tncl(f)&\leq\ncl_{\As(P)}(\relAs(f))\leq\tncl(f)+1.
\end{align*}
If $f$ is not an isomorphism, then $\covlen(f)=\tncl(f)+1$ and hence $\ncl(f) \leq\covlen(f) \leq\ncl(f)+1$.
If $f$ is connected, then $\covlen(f)=\ncl(f)$.
\end{thm}

\begin{proof}
The first inequalities follow from \cref{proposition: nilpotency detected by relative transvection group}, and the second from \cref{proposition: nilpotency detected by relative adjoint group}.
The formula for $\covlen(f)$ follows from \cref{theorem: transvection class and iterated coverings}, and the connected case follows from \cref{corollary: connected nilpotent homomorphism}.
\end{proof}



For a non-isomorphism, the covering length is either $\ncl(f)$ or $\ncl(f)+1$, and the nilpotency class of the relative transvection group is either $\ncl(f)-1$ or $\ncl(f)$. The following examples show that both possibilities occur.


\begin{eg}
\label[eg]{example: basic nilpotent homomorphisms}
Let $f\colon X\twoheadrightarrow Y$ be a non-bijective surjection between trivial quandles. Then $\Inn(X)=\Inn(Y)=1$, so the map is rigid and hence $0$-nilpotent. It is also a covering, but not an isomorphism. Thus $\ncl(f)=\tncl(f)=0$ and $\covlen(f)=1$.
\end{eg}

\begin{eg}
\label[eg]{example: two possible class gaps}
Let
$f\colon\Lambda_{4, 3}\twoheadrightarrow\Lambda_{2,1}$ be reduction
modulo $2$ between linear Alexander quandles. If $x\equiv y\pmod2$, then
$(1-3)(x-y)\equiv0\pmod4$, so $s_x=s_y$. Hence $f$ is a covering and
$\relTrans(f)=1$. The target is trivial, while
$\Inn(\Lambda_{4,3})$ is a nontrivial abelian group. Therefore
$\ncl(f)=\covlen(f)=1$ and $\tncl(f)=0$.

%Together with the non-bijective surjection between trivial quandles in \cref{example: basic nilpotent homomorphisms}, this realizes both possible gaps in \cref{theorem: main structure theorem for nilpotent homomorphisms}.

Note that, although the only connected nilpotent quandle is the one-point trivial quandle $\{\ast\}$ (\cite[Corollary 3.7]{darne_2026_nilpotent_quandles}), 
this homomorphism is connected and $1$-nilpotent but not an isomorphism; compare \cref{proposition: strongly connected and nilpotent is iso}.
\end{eg}



We next consider the nilpotency under base change. 
Recall that covering homomorphisms are stable under arbitrary pullbacks by \cref{proposition: pull-back of covering is again covering}, and rigid homomorphisms are stable under pullbacks along surjections by \cref{theorem: connected-rigid factorization}.

\begin{prop}
\label[prop]{proposition: nilpotency and covering length under pullback}
Let $f\colon P\twoheadrightarrow Q$ be a nilpotent homomorphism, $v\colon R\to Q$ a quandle homomorphism, and $f'\colon P\times_QR\twoheadrightarrow R$ the pullback of $f$ along $v$.
Then the following hold.
\begin{enumerate}
    \item\label{item: pullback general inequalities} $\covlen(f')\leq\covlen(f)$, and $\ncl(f')\leq \ncl(f)+1$.
    \item\label{item: pullback surjective case} If $v$ is surjective, then $\covlen(f')=\covlen(f)$, $\tncl(f')=\tncl(f)$, and $\ncl(f')\leq \ncl(f)\leq \ncl(f')+1$.
    \item\label{item: pullback connected case} If $f$ is connected, $\ncl(f')\leq\ncl(f)$.
\end{enumerate}
\end{prop}

\begin{proof}
\eqref{item: pullback general inequalities}: Pulling back a covering factorization gives $\covlen(f')\leq\covlen(f)$ by \cref{proposition: pull-back of covering is again covering}. Hence, $\ncl(f')\leq\covlen(f')\leq\covlen(f)\leq\ncl(f)+1$ by \cref{theorem: comparison of ncl tncl covlen}.

\eqref{item: pullback surjective case}: Let $u\colon P\times_Q R\twoheadrightarrow P$ be the projection. By \cite[Proposition~3.30]{preprint_yuki_tomoki_2026_relativization_of_symmetries_on_quandles}, $u_*$ restricts to an isomorphism $\relTrans(f')\xrightarrow{\sim}\relTrans(f)$. \cref{proposition: functoriality of relative lower central series} and this give $u_*\colon\relTransLCS{i}{f'}\xrightarrow{\sim}\relTransLCS{i}{f}$ for every $i\geq0$. Hence, $\tncl(f')=\tncl(f)$. Since $v$ is surjective, $f'$ is an isomorphism if and only if $f$ is. If they are isomorphisms, both covering lengths are $0$. Otherwise, $\covlen(f')=\tncl(f')+1=\tncl(f)+1=\covlen(f)$ by \cref{theorem: transvection class and iterated coverings}.

The inequality $\ncl(f')\leq\ncl(f)$ follows by pulling back the factorization in \cref{theorem: nilpotent homomorphism and covering rigid tower}. By \cref{theorem: comparison of ncl tncl covlen}, we have $\ncl(f)\leq \tncl(f)+1=\tncl(f')+1\leq\ncl(f')+1$.

\eqref{item: pullback connected case}: If $f$ is connected, then $\covlen(f)=\ncl(f)$ by \cref{theorem: comparison of ncl tncl covlen}. Therefore, $\ncl(f')\leq\covlen(f')\leq\covlen(f)=\ncl(f)$.
\end{proof}




\begin{eg}
\label[eg]{example: sharpness of pullback class estimate}
Let $G=Q_8\rtimes\langle\sigma\rangle$, where $Q_8$ is the quaternion group and $\sigma$  cyclically permutes $i,j,k\in Q_8$ and has order $3$. Then, $Z(G)=\{\pm1\}$ and $H\coloneqq G/Z(G)\cong A_4$. 
The quotient $\pi\colon G\twoheadrightarrow H$ induces $f\coloneqq \Conj(\pi)\colon\Conj(G)\twoheadrightarrow\Conj(H)$. 
Since $Z(H)=1$, we have $\relInn(f)=1$, and thus, $f$ is rigid and $\ncl(f)=0$.

Let $V=Q_8/Z(G)\cong V_4\subseteq H$. Pulling back $f$ along $\Conj(V)\hookrightarrow\Conj(H)$ gives $f'\colon \Conj(Q_8)\twoheadrightarrow\Conj(V_4)$. 
As $\Conj(V_4)$ is a trivial quandle, 
\[
\relInn(f') = \Inn(\Conj(Q_8)) \cong Q_8/Z(Q_8)\cong V_4.
\] 
Hence the pullback is $1$-nilpotent but not rigid. 
It shows that the inequality in \eqref{item: pullback general inequalities} of \cref{proposition: nilpotency and covering length under pullback} is sharp.
\end{eg}

We end this section with an explicit computation for Alexander homomorphisms and an application showing that the upper bound for the nilpotency class under composition is also sharp.

\begin{prop}
\label[prop]{example: relative Alexander lower central series}
    In the setting in \cref{example: Alexander quandle}, put $f\coloneqq\Alex(\Phi)$ and $D\coloneqq(1-\sigma)A\cap K$. Then, we have $\relInnLCS{i}{f} = \{L_d\mid d\in(1-\sigma)^iD\}$ for $i\geq1$ and $\relTransLCS{i}{f} = \{L_e\mid e\in(1-\sigma)^{i+1}K\}$ for $i\geq0$.
\end{prop}
\begin{proof}
Put $t\coloneqq1-\sigma$, $G\coloneqq\Inn(\Alex(A,\sigma))$,
and $L(U)\coloneqq\{L_u\mid u\in U\}$.
Since $G$ is generated by $\sigma$ and $L(tA)$, the translations
commute, and $[\sigma,L_u]=L_{-tu}$, we have $[G,L(U)]=L(tU)$
for every $\sigma$-invariant subgroup $U\subseteq tA$.

For $x\in A$ and $\varphi\in\relInn(f)$, we have
$\varphi(x)-x\in D$, since $\varphi$ preserves the cosets of $tA$
and the fibers of $f$. Thus,
$[\varphi,s_x]=s_{\varphi(x)}s_x^{-1}
=L_{t(\varphi(x)-x)}\in L(tD)$.
Since $L(tD)$ is normal in $G$, this gives
$[G,\relInn(f)]\subseteq L(tD)$.
The reverse inclusion follows from $L(D)\subseteq\relInn(f)$
and $[G,L(D)]=L(tD)$.
The two formulas now follow by induction, using
$\relTrans(f)=L(tK)$ from \cref{example: Alexander quandle}.
\end{proof}


From \cref{proposition: nilpotency under composition}, $\ncl(g\circ f)\leq \ncl(f)+\ncl(g)$ in general. The following shows that the equality can occur.

\begin{eg}
\label[eg]{example: sharpness of composition class estimate}
Let $m, n$ be integers such that $m>n\ge2$ and $R_n\coloneqq \Lambda_{n, -1}$ be the dihedral quandle of order $n$.
Consider the quandle homomorphism $h_{m, n}\colon R_{2^m} \twoheadrightarrow R_{2^n}$ induced by the reduction group homomorphism $\Phi\colon \ZZ/2^m\ZZ\twoheadrightarrow\ZZ/2^n\ZZ$.
Since $\Ker(\Phi) = 2^n\ZZ/2^m\ZZ$, $D= 2(\ZZ/2^m\ZZ)\cap \Ker(\Phi) = 2^n\ZZ/2^m\ZZ$. 
Therefore, we have $\relInnLCS{i}{h_{m, n}} = \{L_x \mid x\in 2^{n+i}\ZZ/2^m\ZZ\}$ by \cref{example: relative Alexander lower central series}. 
Thus, we have $\ncl(h_{m, n}) = m-n$.
Applying it to the composition $R_{2^{m+n+2}}\xrightarrow{f} R_{2^{m+2}}\xrightarrow{g} R_{4}$, we have
$\ncl(f)=n$, $\ncl(g)=m$, and $\ncl(g\circ f)=m+n$. Thus, $\ncl(g\circ f)=\ncl(f)+\ncl(g)$ holds.
\end{eg}

